% =========================================================
% Formula Depth (Complexity)
% =========================================================

\section{Syntax: Formula Depth}

% ---------------------------------------------------------
% TOOLKIT
% ---------------------------------------------------------

% ---------------------------------------------------------
% Formula Depth
% ---------------------------------------------------------
\begin{definitionbox}{Definition (Formula Depth)}
\begin{definition}[Formula Depth]\label{def:formula-depth}
The \emph{depth} (or \emph{complexity}) of a formula $\varphi$, denoted
$\mathrm{depth}(\varphi)$, is a natural number defined recursively:
\begin{enumerate}
  \item \textbf{Atomic case.} If $\varphi$ is atomic, then $\mathrm{depth}(\varphi) = 0$.

  \item \textbf{Negation.} If $\varphi = \neg \psi$, then
  $\mathrm{depth}(\varphi) = \mathrm{depth}(\psi) + 1$.

  \item \textbf{Binary connectives.} If $\varphi = (\psi \circ \chi)$
  for $\circ \in \{\wedge,\vee,\rightarrow,\leftrightarrow\}$, then
  \[
  \mathrm{depth}(\varphi) = \max\{\mathrm{depth}(\psi),\,\mathrm{depth}(\chi)\} + 1.
  \]

  \item \textbf{Quantifiers.} If $\varphi = \forall x\,\psi$ or $\varphi = \exists x\,\psi$,
  then $\mathrm{depth}(\varphi) = \mathrm{depth}(\psi) + 1$.
\end{enumerate}
\end{definition}
\end{definitionbox}

\begin{dependencies}
\begin{itemize}
  \item TODO: determine dependencies for \texttt{def:formula-depth}.
\end{itemize}
\end{dependencies}

\begin{remark*}[English reading]
The depth of a formula measures the maximum number of logical formation steps
required to build the formula from atomic formulas. It corresponds to the height
of the formula's syntactic parse tree.
\end{remark*}

\begin{remark*}[Proof strategy]
Formula depth is the standard measure used in structural induction proofs over
first-order formulas. The base case is depth~$0$ (atomic formulas); the
inductive step handles each formation rule and assumes the result holds for
sub-formulas of strictly smaller depth.
\end{remark*}
